\begin{proposition}[the logarithmic Kodaira--Spencer divisor]
\label{prop:log-ks-divisor}
\Verified\ The connecting maps of (11.2) are
\begin{align*}
 \kappa_0&:f_*\mathcal E\longrightarrow
  \mathcal L\otimes R^1f_*\mathcal O_X,\\
 \kappa_1&:R^1f_*\mathcal E\longrightarrow
  \mathcal L\otimes R^2f_*\mathcal O_X\simeq\mathcal O.
\end{align*}
Then
\[
 \operatorname{div}(\det\kappa_0)=p_1+p_2,
 \qquad \kappa_1\text{ is surjective as a sheaf}.
\tag{11.6}
\]
In particular $\kappa_0$ is injective and
$\operatorname{coker}\kappa_0=\mathbb C_{p_1}\oplus\mathbb C_{p_2}$.
Moreover the raw period determinant has
\[
 \operatorname{ord}_{z_1}K=2,\qquad
 \operatorname{ord}_{z_2}K=0.
\tag{11.7}
\]
\end{proposition}

\begin{proof}
On the smooth locus, the connecting construction and its
exterior-product Leibniz rule are the Gauss--Manin/Kodaira--Spencer
formalism of
\cite[\S2, equations~(5)--(10), pp.~202--204]{KatzOda1968}.
Katz--Oda work with a smooth algebraic morphism; the citation is not being
used as a logarithmic analytic special-fibre theorem.  In the present family,
Deligne's local monomial calculation supplies the locally split logarithmic
sequence (11.2) \cite[Chapter~II, \S3.3]{Deligne1970}.  Taking a local
splitting and a \v Cech representative, the connecting cochain is the
difference of the two splittings.  Since exterior product of cochains obeys
\[
 \delta(\alpha\wedge\beta)
 =\delta\alpha\wedge\beta+(-1)^{\deg\alpha}
   \alpha\wedge\delta\beta,
\]
the two displayed connecting maps extend across every logarithmic chart and
satisfy the same Leibniz rule.  Thus the singular-fibre formula below is a
direct calculation from (11.2), not an extrapolation of the Katz--Oda
hypotheses.
The determinant is a section of
\[
 \det(\mathcal L\otimes G)\otimes\det(f_*\mathcal E)^{-1}
 =\mathcal L^{\otimes2}\otimes\det G
   \otimes\det(f_*\mathcal E)^{-1}
 \simeq\mathcal O(2)\otimes\mathcal O(-1)\otimes\mathcal O(1)
 \simeq\mathcal O(2).
\tag{11.8}
\]
At the cusp, $Z=sB_0+C(t)$ and $dZ/ds=B_0+O(t)$, so it is a unit.  At a
finite point use $\ell=dt/t=m\,ds/s$.  In invariant determinant frames
with no normalization imposed against the raw marked frames, direct
calculation gives
\[
 \det\kappa_0
 =\varepsilon(t)m^{-2}s^{m+2-2a}K(s),
 \qquad \varepsilon(0)\ne0.
\tag{11.9}
\]
Here $\varepsilon$ is the nowhere-zero invariant holomorphic unit recording the two
frame normalizations; compatible eigenframes may be normalized so that
$\varepsilon=1$.  The dot in $K$ means $d/ds$ in the chosen linearizing coordinate.
Replacing $s$ by another such coordinate changes only the unit, so every
divisor and order calculation below is coordinate-independent.
Equivariance gives
\[
 \operatorname{ord}_sK\equiv2a-2\pmod m.
\tag{11.10}
\]
At $(m,a)=(3,2)$, 
\cref{prop:period-derivative-determinant} supplies $K(z_1)=0$; hence the
right side of (11.9) has coarse order at least one.  At $(m,a)=(4,1)$,
the prefactor is $s^4=t$, so the coarse order is again at least one.
The two forced zeros exhaust the degree in (11.8).  They are therefore
simple, there are no others, and (11.7) follows from (11.9)--(11.10).

For completeness, we derive the coarse matrices rather than only recording
their shape.  Let
\[
 D_{\mathrm{KS}}(s)=(d_{rs}(s))_{0\leq r,s\leq1}
\]
be the Kodaira--Spencer matrix on the marked cover after the common factor
$ds$ has been removed.  Since $t=s^m$ and
$ds=(s/m)\ell$, passage from the marked frames to the actual invariant
frames in (11.4) gives
\[
 A=\frac{s}{m}d_{00},\qquad
 B=\frac{s^{1-a}}m d_{01},\qquad
 C=\frac{s^{1-a}}m d_{10},\qquad
 D_0=\frac{s^{m+1-2a}}m d_{11}.
\tag{11.10a}
\]
The character of each expression in (11.10a) is trivial, so each is a
holomorphic function of $t$.  In the bases
$(e_0,s^{m-a}e_1)$ and $(b_0,s^ab_1)$ one consequently has the exact matrix
\[
 [\kappa_0]=\begin{pmatrix}A&tB\\ C&D_0\end{pmatrix}.
\tag{11.10b}
\]
Its determinant is (11.9).  Applying the exterior-product Leibniz rule to
the four ordered generators
\[
 e_0b_0,\quad s^ae_0b_1,\quad
 s^{m-a}e_1b_0,\quad e_1b_1
\]
now gives, without a change of basis,
\[
 [\kappa_1]=(-C,\ A,\ -D_0,\ B)
\tag{11.11}
\]
into the invariant generator $s^a(b_0\wedge b_1)\ell$.  If (11.11) were
not surjective at a finite point, all four entries would lie in $(t)$ and
$\det\kappa_0$ would have order at least two, contrary to (11.6).  Away
from $p_1,p_2$, the determinant divisor (11.6) makes $\kappa_0$
invertible, so the exterior-product row (11.11) contains a unit and is
surjective there as well.  At the cusp this reduction is not an additional
choice.  Use the logarithmic base
frame
\(\ell=ds=(2\pi i)^{-1}dt/t\), the relative cotangent frame
\((e_0,e_1)=([d\log x_1],[d\log x_2])\) of (11.3), and the quotient frame
\((\gamma,u)\) of \(G\), which extends across the cusp by
\cref{thm:structure-direct-images-audit}.  In these frames the smooth-locus Kodaira--Spencer matrix is the
transpose of the derivative of the period matrix: for
\(h=a e_0+b e_1\),
\[
 \kappa_0(h)=
 (6a\dot\mu+b\dot\beta)\gamma+
 (a\dot\tau+b\dot\mu)u .
 \tag{11.11a}
\]
Here the dot is \(d/ds\).  The cusp expansion
\(Z=sB_0+C(t)\), \(t=e^{2\pi i s}\), gives
\[
 \frac{dZ}{ds}=B_0+2\pi i\,tC'(t),
 \qquad
 B_0=\begin{pmatrix}0&1\\-1&0\end{pmatrix}.
\]
These logarithmic frames extend over \(t=0\), so (11.11a) reduces exactly
to \(B_0^{\mathsf t}\).  Write the cusp matrix, separately from the
finite-fibre normalization above, as
\([\kappa_0]_{p_0}=\left(\begin{smallmatrix}A&E\\C&D\end{smallmatrix}\right)\).
In the same ordered exterior basis used for (11.11), its Leibniz row is
\((-C,A,-D,E)\).  Substituting
\((A,E,C,D)=(0,-1,1,0)\) gives the second reduction.  Thus, in the named
frames,
\[
 [\kappa_0]\equiv\begin{pmatrix}0&-1\\1&0\end{pmatrix}\pmod t,
 \qquad [\kappa_1]\equiv(-1,0,0,-1)\pmod t,
\]
exactly.  Under arbitrary holomorphic changes of source and target frames
the matrices change by invertible left and right factors; in particular
invertibility of \(\kappa_0\) and the existence of a unit entry in the row
of \(\kappa_1\) are unchanged.  Thus $\kappa_1$ is surjective everywhere.
\end{proof}
